\section{Критерії прийняття}
\label{sec:kryterii-pryiniattia}

\begin{longtable}{L{2cm} L{2.5cm} L{9.5cm}}
\caption{Критерії прийняття вимог} \label{tab:ac} \\
\toprule
\textbf{AC ID} & \textbf{Вимога} & \textbf{Критерій прийняття} \\
\midrule
\endfirsthead
\toprule
\textbf{AC ID} & \textbf{Вимога} & \textbf{Критерій прийняття} \\
\midrule
\endhead
\midrule
\multicolumn{3}{r}{\emph{Продовження на наступній сторінці}} \\
\endfoot
\bottomrule
\endlastfoot

AC-001 & REQ-F-001 & Після заповнення коректних email, пароля та нікнейму і підтвердження реєстрації користувач отримує доступ до порожнього особистого каталогу. Повторна реєстрація з тим самим email неможлива: система відображає повідомлення про існуючий акаунт. \sked{K7} \\
\addlinespace
AC-002 & REQ-F-002 & Після введення правильних email і пароля користувач перенаправляється до свого каталогу. При неправильних даних~--- повідомлення про помилку, вхід не виконується. \sked{K8} \\
\addlinespace
AC-003 & REQ-F-003 & Після виходу звернення до захищених сторінок перенаправляє на сторінку входу. Повторна авторизація відновлює доступ. \sked{K9} \\
\addlinespace
AC-004 & REQ-F-004 & Каталог відображає всі додані елементи; для кожного видно назву, тип, статус і оцінку. \sked{K10} \\
\addlinespace
AC-005 & REQ-F-005 & При виборі типу «Книги» відображаються лише елементи типу «книга»; при скиданні фільтра~--- знову всі. \sked{K11} \\
\addlinespace
AC-006 & REQ-F-006 & При виборі статусу «завершено» відображаються лише елементи з цим статусом; при скиданні~--- знову всі. \sked{K12} \\
\addlinespace
AC-007 & REQ-F-007 & Після заповнення обов'язкових полів (назва, тип) елемент з'являється в каталозі з відповідними атрибутами та обкладинкою. \sked{K13} \\
\addlinespace
AC-008 & REQ-F-008 & Після редагування та збереження змінені атрибути відображаються в каталозі й на сторінці елемента. \sked{K14} \\
\addlinespace
AC-009 & REQ-F-009 & Після підтвердження видалення елемент зникає з каталогу та з усіх добірок, до яких входив. \sked{K15} \\
\addlinespace
AC-010 & REQ-F-010 & Доданий тег відображається на картці. При введенні в поле тегу~--- підказки раніше використаних тегів. Видалення тегу не впливає на інші атрибути. \sked{K16} \\
\addlinespace
AC-011 & REQ-F-011 & При фільтруванні за тегом відображаються лише позначені ним елементи; результат оновлюється при зміні або скиданні фільтра. \sked{K17} \\
\addlinespace
AC-012 & REQ-F-012 & Для елемента з двома і більше Sessions~--- список усіх записів з датою, оцінкою та нотаткою. \sked{K18} \\
\addlinespace
AC-013 & REQ-F-013 & Sessions відсортовані від найновішої до найстарішої; загальна кількість вказана. \sked{K19} \\
\addlinespace
AC-014 & REQ-F-014 & Після редагування Session~--- оновлені дані. Після видалення~--- сесія зникає, кількість зменшується. \sked{K20} \\
\addlinespace
AC-015 & REQ-F-015 & Елемент одночасно відображається у кожній добірці, до якої його додано. Видалення з однієї не впливає на інші. \sked{K21} \\
\addlinespace
AC-016 & REQ-F-016 & Після введення назви нова порожня добірка з'являється в списку добірок. \sked{K22} \\
\addlinespace
AC-017 & REQ-F-017 & Після збереження нова назва відображається в списку і в заголовку сторінки. \sked{K23} \\
\addlinespace
AC-018 & REQ-F-018 & Після підтвердження видалення добірка зникає; елементи залишаються в каталозі. \sked{K24} \\
\addlinespace
AC-019 & REQ-F-019 & Після видалення елемента з добірки він залишається в каталозі та в інших добірках. \sked{K25} \\
\addlinespace
AC-020 & REQ-F-020 & Сторінка добірки відображає всі елементи з назвою, типом і статусом. \sked{K26} \\
\addlinespace
AC-021 & REQ-F-021 & Якщо дату не заповнено при зміні статусу на «завершено»~--- автоматично підставляється поточна дата. Якщо введено вручну~--- зберігається введене. \sked{K27} \\
\addlinespace
AC-022 & REQ-F-022 & Сторінка статистики відображає кількість завершених елементів за типами та за обраний проміжок. \sked{K28} \\
\addlinespace
AC-023 & REQ-F-023 & При введенні частини назви каталог відображає лише відповідні елементи; при очищенні~--- всі. \sked{K29} \\
\addlinespace
AC-024 & REQ-F-024 & Форма відображає тип-специфічні поля лише після вибору типу; усі поля необов'язкові. \sked{K30} \\
\addlinespace
AC-025 & REQ-F-025 & Без власного зображення~--- дефолтна ілюстрація за типом. Після завантаження~--- власне зображення; його можна видалити, повернувши дефолтне. \sked{K31} \\
\addlinespace
AC-026 & REQ-F-026 & Після вибору параметра сортування каталог перебудовується; при зміні~--- знову. \sked{K32} \\
\addlinespace
AC-027 & REQ-F-027 & Після вибору параметра сортування всередині добірки список перебудовується. \sked{K33} \\
\addlinespace
AC-028 & REQ-F-028 & Спроба завантажити файл невідповідного формату або розміром понад 5~МБ~--- повідомлення про помилку, файл не зберігається. \sked{K48} \\
\addlinespace
AC-029 & REQ-F-029 & Після зміни нікнейму в налаштуваннях акаунту нове ім'я відображається в інтерфейсі. \sked{K49} \\
\addlinespace
AC-030 & REQ-F-022a & Сторінка статистики містить перемикач часового проміжку: «цей місяць», «цей рік», «весь час». \sked{K51} \\
\addlinespace
AC-031 & REQ-NF-006 & При введенні пароля коротше 8~символів~--- конкретне повідомлення. При некоректному email або нікнеймі поза межами 2–50~символів~--- відповідне повідомлення. \sked{K46} \\
\addlinespace
AC-032 & REQ-NF-008 & При порожньому каталозі, порожній добірці або відсутніх результатах пошуку відображається інформативне повідомлення (не порожній список). \sked{K52} \\
\addlinespace
AC-033 & REQ-NF-009 & Після запуску тестів у директорії \texttt{logs/} існує файл з результатами останнього запуску. Файл містить: загальну кількість тестів, кількість passed~/ failed~/ skipped, статус кожного тесту (назва~+ результат), час виконання набору. \sked{K54} \\

\end{longtable}
